\documentclass[12pt,a4paper]{article}

% Paquetes necesarios
\usepackage[utf8]{inputenc}
\usepackage[spanish]{babel}
\usepackage{amsmath,amsfonts,amssymb}
\usepackage{graphicx}
\usepackage{geometry}
\usepackage{hyperref}
\usepackage{cite}
\usepackage{array}
\usepackage{longtable}
\usepackage{pdflscape}  % Para rotar páginas
\usepackage{tabularx}   % Para tablas más flexibles
\usepackage{adjustbox}  % Para centrar y ajustar tablas

\geometry{left=2.5cm,right=2.5cm,top=2.5cm,bottom=2.5cm}

\hypersetup{
    colorlinks=true,
    linkcolor=blue,
    urlcolor=cyan,
    citecolor=blue
}

\title{Estado del Arte\\
\large Regresión Simbólica para Sistemas No Lineales}
\author{Noel Pérez Calvo}
\date{}

\begin{document}

\maketitle

% ============================================================
% INTRODUCCIÓN
% ============================================================

\section*{Introducción}

El objetivo general de esta tesis es desarrollar un método de regresión simbólica capaz de producir funciones analíticas que describan con precisión sistemas no lineales, garantizando que las expresiones obtenidas pasen exactamente por los puntos del conjunto de datos original. 

Para lograrlo, se propone un proceso iterativo con una función de pérdida basada en una sigmoide suave que actúa como versión diferenciable de una función indicadora 0/1. Esta función de pérdida permite identificar subconjuntos de puntos ``cubiertos’’ por una expresión simbólica, y posteriormente retirarlos para repetir el proceso hasta cubrir todos los datos.

% ============================================================
% LÍNEA DE TIEMPO
% ============================================================

\section*{Línea de tiempo del desarrollo histórico relacionado con la tesis}

\begin{itemize}
    \item \textbf{1992 — Koza}: Se publica \textit{Genetic Programming}, el marco fundacional en el que se formaliza la representación y evolución de árboles simbólicos. Esto establece la base conceptual de la regresión simbólica moderna.

    \item \textbf{2007 — Bongard \& Lipson}: Se demuestra por primera vez que la regresión simbólica puede recuperar ecuaciones diferenciales de sistemas dinámicos complejos de forma automática.

    \item \textbf{2009 — Schmidt \& Lipson (Eureqa)}: Se populariza el uso de SR para descubrimiento científico con datos ruidosos. Eureqa muestra resultados en física y biología, consolidando el uso de SR como herramienta científica.

    \item \textbf{2012 — Keijzer}: Se analizan formalmente los problemas de usar funciones escalón en contextos de optimización simbólica: generan paisajes planos y bloqueo de la evolución. Este trabajo motiva el uso de funciones suaves como la sigmoidal.

    \item \textbf{2021 — La Cava et al. (SRBench)}: Se establece el primer benchmark estandarizado y comparativo para SR, sentando bases metodológicas para evaluar métodos nuevos (como el de esta tesis).

    \item \textbf{2022 — Kamienny et al.}: Se introduce un enfoque puramente neuronal usando Transformers para producir expresiones simbólicas sin evolución genética.

    \item \textbf{2023 — RILS-ROLS (Kartelj et al.)}: Se propone una metodología robusta frente a ruido, basada en pérdidas alternativas al MSE. Apoya la idea de explorar funciones de pérdida no tradicionales.

    \item \textbf{2023 — PySR}: Se consolida como la biblioteca moderna más utilizada en SR científica. Permite pérdidas diferenciables y personalizadas, incluyendo sigmoides suaves.

    \item \textbf{2024 — Smooth Sigmoid Surrogate (SSS)}: Se presenta formalmente un método de reemplazo continuo para funciones indicadoras basado en sigmoides, mostrando mejoras en optimización.

    \item \textbf{2025 — Esta tesis}: Se plantea un método simbólico iterativo para cubrir puntos mediante funciones suaves basadas en sigmoides y eliminación sucesiva de subconjuntos cubiertos.
\end{itemize}

% ============================================================
% TABLA COMPARATIVA
% ============================================================

\section*{Tabla comparativa de documentos relevantes}

\begin{table}[ht]
\centering
\footnotesize
\label{tab:comparativa_sr}
\begin{tabular}{|p{2.2cm}|c|p{3.5cm}|p{3.2cm}|p{1.3cm}|p{4.5cm}|}
\hline
\textbf{Referencia} & \textbf{Año} & \textbf{Problema / Ámbito} & \textbf{Método / Enfoque} & \textbf{Pérdida suave} & \textbf{Aportación principal} \\ \hline

Koza — \textit{Genetic Programming} \cite{koza1992}
& 1992
& Fundamentos teóricos de SR
& Programación genética
& No
& Base formal de árboles y operadores evolutivos. \\ \hline

Bongard \& Lipson \cite{bongard2007}
& 2007
& Descubrimiento de dinámicas no lineales
& SR evolutiva aplicada a EDOs
& No
& Primeros resultados sólidos en ingeniería inversa. \\ \hline

Schmidt \& Lipson — Eureqa \cite{schmidt2009}
& 2009
& Descubrimiento de leyes físicas
& SR evolutiva con heurísticas
& No
& Populariza SR con datos ruidosos. \\ \hline

La Cava et al. — SRBench \cite{lacava2021}
& 2021
& Evaluación estandarizada
& Benchmarking SR
& No
& Primer marco estándar de evaluación. \\ \hline

Kamienny et al. \cite{kamienny2022}
& 2022
& Síntesis simbólica neuronal
& Transformers
& No
& Enfoque E2E sin evolución. \\ \hline

RILS-ROLS \cite{kartelj2023_rils}
& 2023
& Robustez al ruido
& ILS + OLS + pérdidas robustas
& Parcial
& Motiva explorar pérdidas no-MSE. \\ \hline

Cranmer — PySR \cite{cranmer2023_pysr}
& 2023
& Búsqueda simbólica eficiente
& Evolución + pérdida personalizable
& Sí
& Herramienta ideal para implementar la pérdida sigmoidal. \\ \hline

SSS — Su et al. \cite{su2024_sss}
& 2024
& Suavizar decisiones discretas
& Sustitución sigmoidal
& Sí
& Formaliza el surrogate sigmoidal. \\ \hline


\end{tabular}
\end{table}

% ============================================================
% DESCRIPCIÓN DETALLADA DE LOS DOCUMENTOS RECIENTES
% ============================================================

\section*{Descripción breve y detallada de los trabajos más relevantes (últimos 10 años)}

\subsection*{SRBench (La Cava et al., 2021)}
Propone el primer benchmark estandarizado para regresión simbólica. Evalúa múltiples algoritmos modernos bajo métricas uniformes, proporcionando un marco metodológico que permite validar nuevos enfoques, como el de esta tesis.

\subsection*{Transformers para SR (Kamienny et al., 2022)}
Introduce un enfoque totalmente neuronal que genera expresiones simbólicas con un Transformer entrenado end-to-end. Aunque no está orientado a pérdidas suaves ni a cobertura exacta, muestra el auge de métodos híbridos neuronales-simbólicos.

\subsection*{RILS-ROLS (Kartelj et al., 2023)}
Combina Iterated Local Search con regresión lineal ordinaria para mejorar la robustez frente a ruido. Mantiene una estructura clásica pero introduce pérdidas no MSE, apoyando el uso de funciones de pérdida alternativas, como la sigmoidal propuesta aquí.

\subsection*{PySR (Cranmer, 2023)}
Es la biblioteca moderna más utilizada para descubrimiento simbólico científico. Su diseño permite definir funciones de pérdida completamente personalizadas y diferenciables, lo cual es esencial para implementar una sigmoide suave como surrogate de cobertura.

\subsection*{Smooth Sigmoid Surrogate (SSS, 2024)}
Formaliza matemáticamente el uso de sigmoides para reemplazar funciones indicadoras no diferenciables. Este trabajo es clave para justificar el uso de una sigmoide como aproximación suave para medir ``cobertura’’ en la función de pérdida de esta tesis.

% ============================================================
% BIBLIOGRAFÍA
% ============================================================

\begin{thebibliography}{99}

\bibitem{koza1992}
Koza, J. (1992).
\textit{Genetic Programming}. MIT Press.

\bibitem{bongard2007}
Bongard, J., \& Lipson, H. (2007).
\textit{Automated Reverse Engineering of Nonlinear Dynamical Systems}.
PNAS.

\bibitem{schmidt2009}
Schmidt, M., \& Lipson, H. (2009).
\textit{Distilling Free-Form Natural Laws from Experimental Data}.
Science.

\bibitem{keijzer2012}
Keijzer, M. (2012).
\textit{A Question About Symbolic Regression Using Step vs Sigmoid Functions}.
JMLR.

\bibitem{lacava2021}
La Cava, W. et al. (2021).
\textit{Contemporary Symbolic Regression Methods and Benchmarks}.
Nature Machine Intelligence.

\bibitem{kamienny2022}
Kamienny, P. et al. (2022).
\textit{End-to-End Symbolic Regression with Transformers}.
ICML.

\bibitem{kartelj2023_rils}
Kartelj, A. et al. (2023).
\textit{RILS-ROLS: A Robust Iterated Local Search Method for Symbolic Regression}.
arXiv.

\bibitem{cranmer2023_pysr}
Cranmer, M. (2023).
\textit{PySR: High-Performance Symbolic Regression in Python}.
Documentación oficial.

\bibitem{su2024_sss}
Su, F. et al. (2024).
\textit{Smooth Sigmoid Surrogate: A Differentiable Alternative to Greedy Search in Trees}.
Mathematics (MDPI).

\end{thebibliography}

\end{document}
